% ============================================================
%  CHAPITRE : ÉQUATIONS DIFFÉRENTIELLES
%  Mazava Maths — Terminale S (Bacc Série S, Madagascar)
%  Aligné sur le modèle type et le programme officiel MESUPRES
%  (y'=ky, y'+ay=0, y'+ay=b ; ay''+by'+cy=0 : équation
%  caractéristique, 3 cas ; y''=m^2 y, y''=-m^2 y ; conditions
%  initiales ; vérifier qu'une fonction est solution ; second
%  membre : solution particulière + homogène).
%  Environnements : definition, propriete, theoreme, methode,
%                   attention, astucebac, exemple, exercice,
%                   solution, remarque
% ============================================================

\chapter{Équations différentielles}

\begin{citationchapitre}
	« Les lois de la nature ne sont rien d'autre que des relations entre une
	grandeur et ses variations. »\\[0.4em]
	— adage de physique mathématique
\end{citationchapitre}

\begin{objectifs}
	\begin{itemize}[label=\textbullet]
		\item Reconnaître une équation différentielle et son \textbf{ordre}.
		\item Résoudre $y'=ky$, $y'+ay=0$, $y'+ay=b$ ($a$, $b$ constants).
		\item \textbf{Vérifier} qu'une fonction donnée est solution.
		\item Résoudre $ay''+by'+cy=0$ par l'\textbf{équation caractéristique}
		(trois cas) ; cas $y''=m^{2}y$ et $y''=-m^{2}y$.
		\item Déterminer la solution vérifiant des \textbf{conditions initiales}
		(problème de Cauchy).
		\item Traiter une équation \textbf{avec second membre} : solution
		particulière $+$ solution générale de l'équation homogène.
	\end{itemize}
\end{objectifs}


% ================================================================
\section{Généralités}
% ================================================================

\begin{definition}
	Une \textbf{équation différentielle} est une équation dont l'inconnue est
	une \textbf{fonction} $y$ et qui fait intervenir $y$ et ses dérivées
	($y'$, $y''$,\ldots). Son \textbf{ordre} est l'ordre de dérivation le plus
	élevé qui y figure. \textbf{Résoudre} (ou \textbf{intégrer}) l'équation
	sur $I$, c'est en déterminer \emph{toutes} les fonctions solutions sur $I$.
\end{definition}

\begin{methode}
	\textbf{Vérifier qu'une fonction $\varphi$ est solution de $(E)$ :}
	\begin{enumerate}
		\item calculer $\varphi'$ (et $\varphi''$ si besoin) ;
		\item reporter dans le \textbf{premier membre} de $(E)$ ;
		\item simplifier et vérifier l'égalité avec le second membre.
	\end{enumerate}
\end{methode}

\begin{exemple}
	La fonction $\varphi(x)=\e^{-2x}+3$ est-elle solution de $y'+2y=6$ ?

	$\varphi'(x)=-2\e^{-2x}$, donc
	$\varphi'+2\varphi=-2\e^{-2x}+2\e^{-2x}+6=6$. \ Oui, $\varphi$ est solution.
\end{exemple}


% ================================================================
\section{Équations différentielles du premier ordre}
% ================================================================

\begin{theoreme}[Équation homogène $y'+ay=0$]
	Soit $a\in\R$. Les solutions sur $\R$ de $\;y'+ay=0\;$ (c'est-à-dire
	$y'=-ay$) sont les fonctions
	\[
	x\longmapsto C\,\e^{-ax},\qquad C\in\R.
	\]
	En particulier, les solutions de $y'=ky$ sont $x\mapsto C\,\e^{kx}$.
\end{theoreme}

\begin{center}
\begin{tikzpicture}[>=Stealth, scale=1.0]
	\begin{scope}
		\clip (-2.3,-2.5) rectangle (3.4,2.6);
		\foreach \c in {-2,-1,1,2}
			\draw[gray!55,thick] plot[domain=-2.3:3.4,samples=90]
			  (\x,{\c*exp(-\x)});
		\draw[very thick,black] plot[domain=-2.3:3.4,samples=110]
		  (\x,{exp(1-\x)});
	\end{scope}
	\draw[->] (-2.5,0) -- (3.7,0) node[right,font=\small] {$x$};
	\draw[->] (0,-2.7) -- (0,2.9) node[above,font=\small] {$y$};
	\draw[gray!60,dashed] (1,0) node[below,font=\scriptsize] {$x_0$}
	  -- (1,1) -- (0,1) node[left,font=\scriptsize] {$y_0$};
	\filldraw (1,1) circle (1.6pt) node[above right,font=\scriptsize] {$A$};
	\node[font=\scriptsize,align=center] at (2.5,1.9)
	  {solution telle que\\ $y(x_0)=y_0$};
	\node[font=\scriptsize,gray] at (-1.7,-1.9) {$y=C\,\e^{-ax}$};
\end{tikzpicture}
\end{center}

\begin{remarque}
	L'équation $y'+ay=0$ a une \textbf{infinité} de solutions : une par valeur
	de $C$ (le \textbf{faisceau} de courbes ci-dessus). Une \textbf{condition
	initiale} $y(x_0)=y_0$ en sélectionne \textbf{une seule}.
\end{remarque}

\begin{theoreme}[Équation complète $y'+ay=b$]
	Soient $a\neq 0$ et $b$ deux réels. Les solutions sur $\R$ de $\;y'+ay=b\;$
	sont les fonctions
	\[
	x\longmapsto C\,\e^{-ax}+\frac{b}{a},\qquad C\in\R.
	\]
\end{theoreme}

\begin{remarque}
	La solution générale de $y'+ay=b$ est la somme de la solution générale de
	l'équation homogène $y'+ay=0$ et d'une \textbf{solution particulière
	constante} $y_0=\dfrac{b}{a}$.
\end{remarque}

\begin{propriete}{Unicité sous condition initiale}{prop_cauchy_1}
	Pour tout $(x_0\,;y_0)$, il existe une \textbf{unique} solution $f$ de
	$y'+ay=b$ vérifiant $f(x_0)=y_0$.
\end{propriete}

\begin{methode}
	\textbf{Problème de Cauchy} $\bigl(y'+ay=b,\ y(x_0)=y_0\bigr)$ :
	\begin{enumerate}
		\item écrire la solution générale $y(x)=C\e^{-ax}+\dfrac{b}{a}$ ;
		\item imposer $y(x_0)=y_0$, résoudre en $C$ ;
		\item conclure.
	\end{enumerate}
\end{methode}

\exerciceresolu
Résoudre $(E):2y'+3y=6$, puis trouver la solution $f$ telle que $f(0)=-1$.

\solution

On divise par $2$ : $y'+\dfrac32 y=3$, donc $a=\dfrac32$, $b=3$.
Solution générale :
\[
y(x)=C\,\e^{-\frac32 x}+\frac{3}{\frac32}=C\,\e^{-\frac32 x}+2,\qquad C\in\R.
\]
Condition : $f(0)=C+2=-1$, d'où $C=-3$ :
\[
\boxed{f(x)=-3\,\e^{-\frac32 x}+2.}
\]


% ================================================================
\section{Équations du second ordre \texorpdfstring{$ay''+by'+cy=0$}{ay''+by'+cy=0}}
% ================================================================

\begin{definition}
	Soient $a$, $b$, $c$ réels avec $a\neq 0$. L'\textbf{équation
	caractéristique} de $(E):ay''+by'+cy=0$ est l'équation du second degré
	\[
	(E_c):\quad a r^{2}+b r+c=0\qquad (r\in\C).
	\]
\end{definition}

\begin{theoreme}[Résolution selon le discriminant]
	Soit $\Delta=b^{2}-4ac$ le discriminant de $(E_c)$. Les solutions sur $\R$
	de $ay''+by'+cy=0$ sont :
	\begin{center}
	\renewcommand{\arraystretch}{2.0}
	\begin{tabular}{>{\raggedright\arraybackslash}p{3.4cm}
	                >{\raggedright\arraybackslash}p{5.0cm}}
	\hline
	$\Delta>0$ : deux racines réelles $r_1\neq r_2$
	& $y(x)=A\,\e^{r_1 x}+B\,\e^{r_2 x}$ \\
	$\Delta=0$ : racine double $r_0=-\dfrac{b}{2a}$
	& $y(x)=(Ax+B)\,\e^{r_0 x}$ \\
	$\Delta<0$ : racines $\alpha\pm\ii\beta$ ($\beta\neq 0$)
	& $y(x)=\e^{\alpha x}\bigl(A\cos\beta x+B\sin\beta x\bigr)$ \\
	\hline
	\end{tabular}
	\end{center}
	($A$, $B$ décrivent $\R$.)
\end{theoreme}

\begin{attention}
	Dans le cas $\Delta<0$, ne pas confondre le coefficient $\alpha$ de
	l'exponentielle (partie réelle de $r$) et la pulsation $\beta$ dans
	$\cos\beta x$, $\sin\beta x$ (partie imaginaire).
\end{attention}

\begin{propriete}{Cas particuliers du programme}{prop_cas_part}
	Soit $m>0$.
	\begin{itemize}[label=\textbullet]
		\item $y''=m^{2}y$ \ ($y''-m^{2}y=0$) : \
		$y(x)=A\,\e^{mx}+B\,\e^{-mx}$.
		\item $y''=-m^{2}y$ \ ($y''+m^{2}y=0$) : \
		$y(x)=A\cos(mx)+B\sin(mx)$.
	\end{itemize}
\end{propriete}

\begin{astucebac}
	Problème de Cauchy au second ordre : on donne \textbf{deux} conditions
	(souvent $y(0)$ et $y'(0)$). Calculer proprement $y'(x)$ en fonction de
	$A$, $B$ \textbf{avant} de substituer.
\end{astucebac}

\exerciceresolu
Résoudre $(E):y''-4y'+13y=0$ avec $y(0)=2$ et $y'(0)=7$.

\solution

$(E_c):r^{2}-4r+13=0$, $\Delta=16-52=-36<0$ : racines
$r=\dfrac{4\pm 6\ii}{2}=2\pm 3\ii$, donc $\alpha=2$, $\beta=3$.
\[
y(x)=\e^{2x}\bigl(A\cos 3x+B\sin 3x\bigr).
\]
$y'(x)=2\e^{2x}\bigl(A\cos 3x+B\sin 3x\bigr)
+\e^{2x}\bigl(-3A\sin 3x+3B\cos 3x\bigr)$.
\[
y(0)=A=2,\qquad y'(0)=2A+3B=7\ \Rightarrow\ 4+3B=7\ \Rightarrow\ B=1.
\]
\[
\boxed{f(x)=\e^{2x}\bigl(2\cos 3x+\sin 3x\bigr).}
\]


% ================================================================
\section{Équations avec second membre}
% ================================================================

\begin{theoreme}[Structure des solutions]
	Soit $(E)$ une équation différentielle linéaire (ordre $1$ ou $2$) et
	$(E_0)$ l'équation \textbf{homogène} associée (sans second membre).
	Si $y_p$ est \textbf{une} solution particulière de $(E)$, alors
	\[
	f \text{ solution de } (E)
	\iff f-y_p \text{ solution de } (E_0).
	\]
	Donc : \ \textbf{solution générale de $(E)$} $=$ \textbf{solution
	particulière} $+$ \textbf{solution générale de $(E_0)$}.
\end{theoreme}

\begin{methode}
	Pour résoudre $(E)$ avec second membre (sujet guidé) :
	\begin{enumerate}
		\item résoudre l'équation homogène $(E_0)$ ;
		\item déterminer une solution particulière $y_p$ (forme donnée dans
		l'énoncé : on identifie les coefficients) ;
		\item la solution générale de $(E)$ est $y_p+$ (solution générale de $(E_0)$) ;
		\item appliquer les conditions initiales.
	\end{enumerate}
\end{methode}

\exerciceresolu
$(E):y'+2y=(4x+2)\,\e^{-2x}$.
\begin{enumerate}
	\item Résoudre $(E_0):y'+2y=0$.
	\item Déterminer $a$, $b$ tels que $y_p(x)=(a x^{2}+bx)\,\e^{-2x}$ soit
	solution de $(E)$.
	\item En déduire la solution générale de $(E)$, puis la solution $h$
	telle que $h(0)=1$.
\end{enumerate}

\solution

\textbf{1.} $(E_0)$ : $y(x)=C\,\e^{-2x}$, $C\in\R$.

\textbf{2.} $y_p(x)=(a x^{2}+bx)\e^{-2x}$, donc
$y_p'(x)=\bigl[(2ax+b)-2(a x^{2}+bx)\bigr]\e^{-2x}$, et
\[
y_p'+2y_p=(2a x+b)\,\e^{-2x}.
\]
On veut $2a x+b=4x+2$, d'où $a=2$ et $b=2$ :
$y_p(x)=(2x^{2}+2x)\e^{-2x}$.

\textbf{3.} Solution générale : $y(x)=\bigl(2x^{2}+2x+C\bigr)\e^{-2x}$.
$h(0)=C=1$, donc $h(x)=(2x^{2}+2x+1)\,\e^{-2x}$.


% ================================================================
\section{Exercices d'entraînement}
% ================================================================

% --------------------------------------------------
\subsection{Échauffement}
% --------------------------------------------------

\exercice{Lire un faisceau de courbes \hfill \facile}
La figure ci-dessous représente quatre solutions de l'équation
différentielle $(E):y'+y=0$.

\begin{center}
\begin{tikzpicture}[>=Stealth, scale=0.95]
	\begin{scope}
		\clip (-2.7,-2.6) rectangle (2.9,2.6);
		\draw[thick,domain=-1.4:2.7,samples=110] plot (\x,{exp(-\x)});
		\draw[thick,domain=-0.35:2.7,samples=110] plot (\x,{3*exp(-\x)});
		\draw[thick,domain=-2.7:1.4,samples=110] plot (\x,{-exp(-\x)});
	\end{scope}
	\draw[->] (-2.9,0) -- (3.1,0) node[right,font=\small] {$x$};
	\draw[->] (0,-2.8) -- (0,2.9) node[above,font=\small] {$y$};
	\node[font=\scriptsize] at (2.35,0.7) {$C=1$};
	\node[font=\scriptsize] at (1.35,1.55) {$C=3$};
	\node[font=\scriptsize] at (1.9,-0.85) {$C=-1$};
	\node[font=\scriptsize,gray] at (-2.0,0.3) {$C=0$};
\end{tikzpicture}
\end{center}

\begin{enumerate}
	\item Donner la forme générale des solutions de $(E)$.
	\item Identifier la solution qui vérifie $y(0)=3$.
	\item Combien de solutions vérifient $y(0)=0$ ?
	\item Quelle est la limite commune de toutes les solutions en $+\infty$ ?
\end{enumerate}

\exercice{Un faisceau avec second membre \hfill \facile}
La figure ci-dessous représente quatre solutions de l'équation
différentielle $(E):y'+y=2$.

\begin{center}
\begin{tikzpicture}[>=Stealth, scale=0.9]
	\begin{scope}
		\clip (-2.4,-2.4) rectangle (3.4,4.2);
		\foreach \c in {-2,-1,1}
			\draw[thick,domain=-2.4:3.4,samples=110]
			  plot (\x,{2+\c*exp(-\x)});
		\draw[thick,domain=-2.4:3.4,samples=110] plot (\x,{2+2.7*exp(-\x)});
	\end{scope}
	\draw[->] (-2.6,0) -- (3.7,0) node[right,font=\small] {$x$};
	\draw[->] (0,-2.5) -- (0,4.4) node[above,font=\small] {$y$};
	\draw[red!70!black,dashed] (-2.4,2) -- (3.5,2)
	  node[right,font=\scriptsize,red!70!black] {$y=2$};
	\node[font=\scriptsize] at (2.7,2.5) {solutions};
\end{tikzpicture}
\end{center}

\begin{enumerate}
	\item Donner la forme générale des solutions de $(E)$.
	\item Quelle est la limite commune de toutes les solutions en $+\infty$ ?
	Interpréter la droite $y=2$.
	\item Identifier la solution constante.
	\item Déterminer la solution vérifiant $y(0)=0$.
\end{enumerate}

\exercice{Vrai ou faux \hfill \facile}
Répondre par \emph{vrai} ou \emph{faux} en justifiant.
\begin{enumerate}
	\item L'équation $y'=3y$ a pour solutions $x\mapsto C\,\e^{3x}$.
	\item $x\mapsto \e^{2x}$ est solution de $y''-4y=0$.
	\item Les solutions de $y''+4y=0$ sont périodiques.
	\item Une équation différentielle a toujours une infinité de solutions.
	\item Si $y_p$ est une solution de $(E):y'+ay=b$, alors $2y_p$ aussi.
	\item Pour $y''=-9y$, la pulsation des solutions est $9$.
	\item La solution de $y'+2y=0$ avec $y(0)=0$ est la fonction nulle.
\end{enumerate}

% --------------------------------------------------
\subsection{Premier ordre}
% --------------------------------------------------

\exercice{Solution générale \hfill \facile}
Résoudre sur $\R$ (solution générale).
\begin{multicols}{2}
\begin{enumerate}
	\item $y'-7y=0$
	\item $2y'+5y=0$
	\item $y'=-3y$
	\item $y'+3y=6$
	\item $4y'-y=5$
	\item $3y'+2y=0$
\end{enumerate}
\end{multicols}

\exercice{Problème de Cauchy \hfill \moyen}
Déterminer la solution vérifiant la condition indiquée.
\begin{multicols}{2}
\begin{enumerate}
	\item $y'=-2y$, \ $y(0)=3$
	\item $y'-4y=12$, \ $y(0)=-1$
	\item $3y'+2y=4$, \ $y(3)=2$
	\item $y'+y=1$, \ $y(0)=0$
\end{enumerate}
\end{multicols}

\exercice{Décroissance radioactive \hfill \moyen}
La masse $m(t)$ (en g) d'un échantillon radioactif vérifie
$m'(t)=-\lambda\,m(t)$ avec $\lambda>0$, et $m(0)=m_0$.
\begin{enumerate}
	\item Exprimer $m(t)$.
	\item La \emph{demi-vie} $T$ est le temps au bout duquel la masse est
	divisée par $2$. Montrer que $T=\dfrac{\ln 2}{\lambda}$.
	\item Déterminer $\displaystyle\lim_{t\to+\infty}m(t)$.
\end{enumerate}

% --------------------------------------------------
\subsection{Second ordre}
% --------------------------------------------------

\exercice{Équation caractéristique \hfill \moyen}
Résoudre sur $\R$ (solution générale).
\begin{multicols}{2}
\begin{enumerate}
	\item $y''-5y'+6y=0$
	\item $y''+6y'+9y=0$
	\item $y''-2y'+5y=0$
	\item $y''-y=0$
	\item $y''+4y=0$
	\item $y''+2y'+2y=0$
\end{enumerate}
\end{multicols}

\exercice{Cauchy au second ordre \hfill \moyen}
\begin{enumerate}
	\item $y''-3y'+2y=0$, \ $y(0)=4$, \ $y'(0)=0$.
	\item $y''+9y=0$, \ $y(0)=1$, \ $y'\!\left(\dfrac{\pi}{6}\right)=0$.
	Montrer que $f(x)$ s'écrit $A\cos(3x+\varphi)$.
\end{enumerate}

\exercice{Oscillations à Analakely \hfill \moyen}
Un système mécanique oscillant est régi par $(E):y''+16y=0$.
\begin{enumerate}
	\item Donner la forme générale des solutions de $(E)$.
	\item Déterminer la solution $g$ telle que
	$g(0)=\dfrac{\sqrt3}{2}$ et $g'(0)=2$.
	\item Vérifier que $g(t)=\cos\!\left(4t-\dfrac{\pi}{6}\right)$ pour tout
	réel $t$.
\end{enumerate}

% --------------------------------------------------
\subsection{Sujets types Bac}
% --------------------------------------------------

\exercice{Sujet type Bac — équation avec second membre \hfill \difficile}
On considère $(E):y'+2y=(3x+1)\,\e^{-2x}$.

\textbf{Partie A}
\begin{enumerate}
	\item Résoudre $(E_0):y'+2y=0$ sur $\R$.
	\item Déterminer $a$, $b$ tels que $u(x)=(a x^{2}+bx)\,\e^{-2x}$ soit
	solution de $(E)$.
	\item Démontrer que $f$ est solution de $(E)$ si et seulement si $f-u$ est
	solution de $(E_0)$.
	\item En déduire l'ensemble des solutions de $(E)$, puis la solution $h$
	telle que $h(0)=2$.
\end{enumerate}

\textbf{Partie B — Modèle thermique}
La température $T(t)$ (en $^{\circ}$C) d'un composant après $t$ minutes est
\[
T(t)=\left(\frac32 t^{2}+t+25\right)\e^{-2t}+20\qquad (t\geq 0).
\]
\begin{enumerate}
	\setcounter{enumi}{4}
	\item Donner la température initiale.
	\item Déterminer $\displaystyle\lim_{t\to+\infty}T(t)$ et l'interpréter.
	\item Calculer $T'(t)$ et étudier son signe sur $[0\,;+\infty[$.
	\item À quel instant $T$ est-elle maximale ? Donner cette température à
	$0{,}1\,^{\circ}$C près.
\end{enumerate}

\exercice{Sujet type Bac — vérifier une solution, puis résoudre \hfill \difficile}
On considère $(E):y'-2y=\dfrac{-2}{1+\e^{-2x}}$ et la fonction
$g(x)=\e^{2x}\ln\!\bigl(1+\e^{-2x}\bigr)$.

\textbf{Partie A}
\begin{enumerate}
	\item Calculer $g'(x)$ et vérifier que $g$ est solution de $(E)$.
	\item Résoudre l'équation homogène $(E_0):y'-2y=0$.
	\item Montrer que $f$ est solution de $(E)$ si et seulement si $f-g$ est
	solution de $(E_0)$.
	\item En déduire la solution générale de $(E)$.
\end{enumerate}

\textbf{Partie B}
\begin{enumerate}
	\setcounter{enumi}{4}
	\item Déterminer $\displaystyle\lim_{x\to-\infty}g(x)$ et
	$\displaystyle\lim_{x\to+\infty}g(x)$ (poser $X=\e^{-2x}$).
	\item En déduire une asymptote à $\mathcal C_g$.
\end{enumerate}

\exercice{Sujet type Bac — tangente imposée \hfill \difficile}
On considère $(E):y''-3y'+2y=0$.
\begin{enumerate}
	\item Résoudre $(E)$ sur $\R$.
	\item Déterminer la solution $f$ dont la courbe passe par $A(0\,;4)$ et
	admet en $A$ une tangente parallèle à l'axe des abscisses.
	\item Étudier les variations de $f$ sur $\R$ et déterminer
	$\displaystyle\lim_{x\to-\infty}f(x)$, $\displaystyle\lim_{x\to+\infty}f(x)$.
	\item Étudier la convexité de $f$ ; montrer que $\mathcal C_f$ a un point
	d'inflexion.
\end{enumerate}

\exercice{Sujet type Bac — pendule et énergie \hfill \difficile}
Pour de petites oscillations, l'angle $\theta(t)$ d'un pendule vérifie
$(E):\theta''+\omega^{2}\theta=0$ ($\omega>0$).
\begin{enumerate}
	\item Donner la solution générale de $(E)$.
	\item À $t=0$, le pendule est écarté de $\theta_0$ et lâché sans vitesse.
	Déterminer $\theta(t)$.
	\item Montrer que l'accélération $\theta''(t)$ est proportionnelle à
	$\theta(t)$ ; préciser le coefficient.
	\item On pose $E(t)=\bigl(\theta'(t)\bigr)^{2}+\omega^{2}\bigl(\theta(t)\bigr)^{2}$.
	Montrer que $E$ est \textbf{constante}.
\end{enumerate}

\exercice{Sujet type Bac — équation avec second membre trigonométrique \hfill \difficile}
On considère $(E):y''+y=\cos x$.
\begin{enumerate}
	\item Résoudre l'équation homogène $(E_0):y''+y=0$.
	\item Montrer que $u(x)=\dfrac12\,x\sin x$ est une solution particulière
	de $(E)$.
	\item En déduire la solution générale de $(E)$.
	\item Déterminer la solution $f$ telle que $f(0)=0$ et $f'(0)=0$.
	\item Calculer $\displaystyle\int_0^{\pi}f(x)\dd x$.
\end{enumerate}

\exercice{Sujet type Bac — changement d'inconnue \hfill \difficile}
On résout sur $I=]0\,;+\infty[$ l'équation $(E):x^{2}y''+3xy'+y=0$.
On pose $z(t)=y(\e^{t})$ pour tout réel $t$.
\begin{enumerate}
	\item Exprimer $z'(t)$ et $z''(t)$ en fonction de $y'(\e^{t})$,
	$y''(\e^{t})$ et $\e^{t}$.
	\item Montrer que $y$ est solution de $(E)$ sur $I$ si et seulement si $z$
	est solution d'une équation $(E'):z''+2z'+z=0$.
	\item Résoudre $(E')$ sur $\R$.
	\item En déduire les solutions de $(E)$ sur $I$.
\end{enumerate}
